Este trabalho apresentou uma metodologia para a classificação automática dos graus ISUP em imagens histopatológicas de próstata, abordando desafios como desbalanceamento de classes, ruído nos rótulos e a natureza ordinal da progressão tumoral. A função de perda híbrida (Ordinal + Focal) foi o principal fator de melhoria, elevando o $\kappa_{\text{quad}}$ de $0{,}826$ para $0{,}856$ (+3\%) e a acurácia de $0{,}592$ para $0{,}669$ (+7\%) em relação ao \textit{baseline}. Destaca-se que a perda ordinal isolada, combinada à remoção de ruído, já promoveu ganho de 2,5\% no $\kappa_{\text{quad}}$ com alta estabilidade (desvio padrão de $0{,}009$). A EfficientNet-B0 demonstrou o melhor equilíbrio entre desempenho, estabilidade e custo computacional, com os erros concentrados entre classes adjacentes — comportamento coerente com a gradação ordinal do ISUP. Como trabalhos futuros, propõe-se investigar arquiteturas baseadas em \textit{Vision Transformers} e estratégias de \textit{ensemble} para aprimorar a classificação dos graus intermediários e ampliar a generalização clínica do modelo.
